% ================================================================ CHAPTER 2
% OBE 2.1 structure: 2.1 Conceptual Framework, 2.2 Data Collection,
% 2.3 Professional Code of Ethical Issues, 2.4 Economic Decision.
\chapter{Research Methodology}
\label{ch:methodology}

This chapter justifies and explains the methodological approach of the study.
The research follows the Design Science Research Methodology (DSRM): the
problem is identified from the literature (Chapter~\ref{ch:introduction}), an
artifact (the detection framework) is designed and developed, and the artifact
is demonstrated and evaluated against explicit research questions using
established benchmarks. DSRM is appropriate because the contribution of this
thesis is a designed computational artifact whose utility must be demonstrated
empirically, rather than a purely observational study.

\section{Conceptual Framework}
\label{sec:framework}

The life cycle of the project comprises five phases: (1) data acquisition and
preprocessing, (2) generator-free pretraining (QACP), (3) supervised multi-task
training of DAVID-Net, (4) evaluation against baselines and ablations, and
(5) deployment as a public web service. The proposed framework is trained in
two stages: \textbf{Stage~0}, Quadrant-Aware Contrastive Pretraining on
synthetically constructed quadrants; and \textbf{Stage~1}, supervised
multi-task fine-tuning of the full DAVID-Net architecture.

\subsection{Problem Formulation}
Let a clip $x = (x_v, x_a)$ consist of a visual stream $x_v$ and an acoustic
stream $x_a$. The framework predicts:
\begin{align}
  \hat{y}_v &= P(\text{video fake} \mid x) \in [0,1], \\
  \hat{y}_a &= P(\text{audio fake} \mid x) \in [0,1], \\
  \hat{q}   &\in \{\text{RVRA}, \text{RVFA}, \text{FVRA}, \text{FVFA}\}, \\
  \hat{m}_v(t),\ \hat{m}_a(t) &\in [0,1] \quad \text{(per-frame manipulation probability)}.
\end{align}

\subsection{DAVID-Net Architecture}
\label{sec:davidnet}

\begin{figure}[t]
  \centering
  % ------------------------------------------------------------------
% DAVID-Net architecture diagram (TikZ). Input via \input{} inside a
% figure environment. Color code: blue = video branch, orange = audio
% branch, violet = consistency path, green = decision heads.
% ------------------------------------------------------------------
\resizebox{\textwidth}{!}{%
\begin{tikzpicture}[
  font=\scriptsize,
  box/.style={draw, rounded corners=2pt, align=center, inner sep=5pt,
              minimum height=0.95cm, minimum width=1.9cm},
  vbox/.style={box, fill=blue!12},
  abox/.style={box, fill=orange!20},
  cbox/.style={box, fill=violet!14},
  hbox/.style={box, fill=green!16},
  obox/.style={box, fill=gray!10},
  emb/.style={box, fill=white, minimum width=1.1cm},
  arr/.style={-{Stealth[length=2.2mm]}, semithick},
  carr/.style={arr, violet!70!black}
]

% ---------------- video branch (top row) ----------------
\node[vbox] (vin) at (0, 2.4)  {Input video\\(25 fps frames)};
\node[vbox] (vcrop) at (3.0, 2.4) {Face + mouth\\crop \& align};
\node[vbox] (venc) at (6.0, 2.4) {VideoMAE\\encoder (frozen)};
\node[vbox] (vtok) at (9.0, 2.4) {$V \in \mathbb{R}^{L_v \times d}$\\video tokens};

% ---------------- audio branch (bottom row) ----------------
\node[abox] (ain) at (0, -2.4)  {Input audio\\(waveform)};
\node[abox] (aprep) at (3.0, -2.4) {Resample 16 kHz\\+ normalize};
\node[abox] (aenc) at (6.0, -2.4) {WavLM\\encoder (frozen)};
\node[abox] (atok) at (9.0, -2.4) {$A \in \mathbb{R}^{L_a \times d}$\\audio tokens};

% ---------------- missing-modality null tokens ----------------
\node[box, dashed, fill=blue!5, minimum width=1.6cm] (nullv) at (9.0, 4.3)
  {null tokens\\(video absent)};
\node[box, dashed, fill=orange!8, minimum width=1.6cm] (nulla) at (9.0, -4.3)
  {null tokens\\(audio absent)};

% ---------------- sync module (middle) ----------------
\node[cbox] (sync) at (9.0, 0) {Sync module\\(windowed InfoNCE)};

% ---------------- fusion ----------------
\node[box, fill=teal!12, minimum height=4.6cm, minimum width=2.3cm]
  (fusion) at (12.0, 0) {Disentangled\\cross-modal\\fusion\\(self- +\\gated cross-\\attention)};

% ---------------- embeddings ----------------
\node[emb, fill=blue!8]   (zv) at (14.6, 2.4)  {$z_v$};
\node[emb, fill=violet!8] (zc) at (14.6, 0)    {$z_c$};
\node[emb, fill=orange!12](za) at (14.6, -2.4) {$z_a$};

% ---------------- heads ----------------
\node[hbox] (hv)   at (17.2, 3.3)  {$H_v$: video\\real/fake};
\node[hbox] (hq)   at (17.2, 1.1)  {$H_{quad}$:\\4-quadrant};
\node[hbox] (hloc) at (17.2, -1.1) {$H_{loc}$: temporal\\localization};
\node[hbox] (ha)   at (17.2, -3.3) {$H_a$: audio\\real/fake};

% ---------------- outputs ----------------
\node[obox] (ov)   at (20.0, 3.3)  {$P(\text{video fake})$};
\node[obox] (oq)   at (20.0, 1.1)  {RVRA / RVFA /\\FVRA / FVFA};
\node[obox] (oloc) at (20.0, -1.1) {$\hat{m}_v(t),\ \hat{m}_a(t)$\\+ sync curve $s_t$};
\node[obox] (oa)   at (20.0, -3.3) {$P(\text{audio fake})$};

% ---------------- arrows: branches ----------------
\draw[arr] (vin) -- (vcrop);
\draw[arr] (vcrop) -- (venc);
\draw[arr] (venc) -- (vtok);
\draw[arr] (ain) -- (aprep);
\draw[arr] (aprep) -- (aenc);
\draw[arr] (aenc) -- (atok);

% tokens -> sync
\draw[carr] (vtok) -- (sync);
\draw[carr] (atok) -- (sync);

% tokens -> fusion
\draw[arr] (vtok.east) -- (fusion.west |- vtok);
\draw[arr] (atok.east) -- (fusion.west |- atok);
\draw[carr] (sync.east) -- (fusion.west |- sync);

% fusion -> embeddings
\draw[arr]  (fusion.east |- zv) -- (zv);
\draw[carr] (fusion.east |- zc) -- (zc);
\draw[arr]  (fusion.east |- za) -- (za);

% embeddings -> heads
\draw[arr]  (zv) -- (hv.west);
\draw[carr] (zc) -- (hv.west);
\draw[arr]  (za) -- (ha.west);
\draw[carr] (zc) -- (ha.west);
\draw[arr]  (zv) -- (hq.west);
\draw[arr]  (za) -- (hq.west);
\draw[carr] (zc) -- (hq.west);
\draw[arr]  (zv) -- (hloc.west);
\draw[arr]  (za) -- (hloc.west);
\draw[carr] (zc) -- (hloc.west);

% heads -> outputs
\draw[arr] (hv) -- (ov);
\draw[arr] (hq) -- (oq);
\draw[arr] (hloc) -- (oloc);
\draw[arr] (ha) -- (oa);

% null-token substitution arrows
\draw[arr, dashed] (nullv) -- (vtok);
\draw[arr, dashed] (nulla) -- (atok);

% disentanglement + availability-gating annotations
\node[align=center, violet!60!black] at (14.6, -4.4)
  {$\mathcal{L}_{dis}$: $z_c \perp z_v,\ z_c \perp z_a$\\(orthogonality penalty)};
\draw[carr, dashed] (14.6, -4.0) -- (zc.south);
\node[align=center, violet!60!black] at (14.6, 4.3)
  {$z_c \leftarrow 0$ if either\\modality is missing};
\draw[carr, dashed] (14.6, 3.7) -- (zc.north);

\end{tikzpicture}}%

  \caption{Overview of the proposed DAVID-Net architecture. Blue denotes the
  video branch, orange the audio branch, violet the cross-modal consistency
  path, and green the multi-task decision heads. The consistency embedding
  $z_c$ is kept disentangled from the modality-specific authenticity
  embeddings $z_v$ and $z_a$.}
  \label{fig:architecture}
\end{figure}

The overall architecture is shown in Figure~\ref{fig:architecture}.
DAVID-Net is built on the observation that ``is stream $X$ internally
authentic?'' and ``do the two streams agree with each other?'' are different
questions. Fusing them prematurely allows shortcut learning (e.g., labeling
every desynchronized clip fake). The architecture therefore keeps
\emph{modality-specific authenticity} evidence disentangled from
\emph{cross-modal consistency} evidence, combining them only at the decision
heads. This design lets the per-modality heads fire independently---making all
four quadrants reachable, including FVFA clips whose two fake streams are
mutually consistent.

\textbf{Encoders.} The video branch extracts aligned face and mouth crops and
encodes them with a self-supervised VideoMAE backbone into a spatio-temporal
token sequence $V \in \mathbb{R}^{L_v \times d}$. The audio branch encodes the
raw 16\,kHz waveform with WavLM~\cite{chen2022wavlm} into
$A \in \mathbb{R}^{L_a \times d}$. Encoders are largely frozen; only adapters
and the fusion layers are trained heavily.

\textbf{Synchronization--consistency module.} Short overlapping windows of $V$
and $A$ are projected into a shared space and trained with an InfoNCE
contrastive objective in which temporally aligned real windows are positives.
The module outputs a per-window agreement score $s_t$ and a pooled consistency
embedding $z_c$.

\textbf{Disentangled cross-modal fusion.} A cross-modal transformer produces
pooled authenticity embeddings $z_v, z_a$ and the consistency embedding $z_c$.
A disentanglement penalty discourages leakage between the authenticity and
consistency subspaces:
\begin{equation}
  \mathcal{L}_{\text{dis}} = \left|\cos(z_v, z_c)\right| + \left|\cos(z_a, z_c)\right|.
\end{equation}

\textbf{Missing-modality support.} Real-world inputs are frequently
single-modality: voice recordings without video, and muted or music-overlaid
clips without usable speech. An absent stream is replaced by learnable null
tokens, the consistency embedding $z_c$ is zeroed (no cross-modal evidence
exists), and the surviving authenticity head decides alone. During training,
\emph{modality dropout} removes one stream from a sample with probability
$p{=}0.15$ (never both), masking the dropped modality's losses per sample.
This enables audio-only detection (evaluated directly against anti-spoofing
baselines on ASVspoof) and silent-video detection within the same network,
with graceful degradation instead of undefined behavior.

\textbf{Multi-task heads and objective.} Four heads consume the disentangled
representations: video authenticity $H_v([z_v; z_c])$, audio authenticity
$H_a([z_a; z_c])$, quadrant $H_{quad}([z_v; z_a; z_c])$, and per-frame
localization $H_{loc}$. The total Stage-1 objective is
\begin{equation}
  \mathcal{L} = \lambda_1 \mathcal{L}_v + \lambda_2 \mathcal{L}_a
  + \lambda_3 \mathcal{L}_{quad} + \lambda_4 \mathcal{L}_{loc}
  + \lambda_5 \mathcal{L}_{sync} + \lambda_6 \mathcal{L}_{dis},
\end{equation}
where $\mathcal{L}_v, \mathcal{L}_a$ are focal binary cross-entropy losses,
$\mathcal{L}_{quad}$ is cross-entropy, $\mathcal{L}_{loc}$ is frame-level
binary cross-entropy against manipulated-interval masks, and
$\mathcal{L}_{sync}$ is the contrastive synchronization loss applied to genuine
(RVRA) samples.

\subsection{Quadrant-Aware Contrastive Pretraining (QACP)}
\label{sec:qacp}

Detectors trained on generator outputs inherit those generators' fingerprints
and fail on unseen synthesis methods. QACP instead pretrains the embedding
space on forgeries \emph{constructed from pristine data alone}, so the space
encodes what manipulation does to a stream rather than what a particular
generator looks like. Starting from verified-real (RVRA) clips only, five
pseudo-classes are manufactured (Table~\ref{tab:qacp-classes}); the complete
two-stage training flow is illustrated in Figure~\ref{fig:qacp}.

\begin{figure}[t]
  \centering
  % ------------------------------------------------------------------
% QACP two-stage training diagram (TikZ): synthetic quadrant
% construction from pristine clips -> factorized SupCon pretraining
% (Stage 0) -> supervised multi-task fine-tuning (Stage 1).
% ------------------------------------------------------------------
\resizebox{\textwidth}{!}{%
\begin{tikzpicture}[
  font=\scriptsize,
  box/.style={draw, rounded corners=2pt, align=center, inner sep=5pt,
              minimum height=0.9cm, minimum width=2.2cm},
  src/.style={box, fill=gray!12},
  tr/.style={box, fill=yellow!25},
  pc/.style={box, fill=cyan!12, minimum width=2.6cm},
  loss/.style={box, fill=red!10},
  st/.style={box, fill=green!14},
  arr/.style={-{Stealth[length=2.2mm]}, semithick}
]

% source
\node[src, minimum height=1.4cm] (real) at (0, 0)
  {Pristine clips\\(RVRA only)\\\emph{no generator}\\\emph{outputs used}};

% transforms
\node[tr] (cs)  at (4.0, 2.6)  {Vocoder\\copy-synthesis\\(audio)};
\node[tr] (sb)  at (4.0, 0.9)  {Self-blended\\face video};
\node[tr] (bo)  at (4.0, -0.9) {Both\\transforms};
\node[tr] (mm)  at (4.0, -2.6) {Cross-clip real\\audio swap};

% pseudo-classes
\node[pc] (prvra) at (8.0, 4.0)  {pseudo-RVRA\\(untouched)};
\node[pc] (prvfa) at (8.0, 2.4)  {pseudo-RVFA};
\node[pc] (pfvra) at (8.0, 0.8)  {pseudo-FVRA};
\node[pc] (pfvfa) at (8.0, -0.8) {pseudo-FVFA};
\node[pc, fill=violet!14] (pmm) at (8.0, -2.6)
  {pseudo-MISMATCH\\(real + real, desynced)\\$\Rightarrow$ mismatch $\neq$ fake};

% shared encoder + fusion
\node[box, fill=teal!12, minimum height=2.2cm] (net) at (11.6, 0.8)
  {Frozen encoders\\+ fusion\\(cached features)};

% factorized supcon
\node[loss] (lv) at (15.2, 2.6)  {SupCon($z_v$;\\video-auth)};
\node[loss] (la) at (15.2, 0.8)  {SupCon($z_a$;\\audio-auth)};
\node[loss] (lc) at (15.2, -1.0) {SupCon($z_c$;\\sync-state)};

% stages
\node[st, minimum height=1.2cm] (s0) at (18.6, 0.8)
  {$\mathcal{L}_{QACP}$\\\textbf{Stage 0}\\pretraining};
\node[st, minimum height=1.2cm] (s1) at (18.6, -2.0)
  {\textbf{Stage 1}\\supervised\\multi-task\\fine-tuning};

% arrows
\draw[arr] (real.east) -- (cs.west);
\draw[arr] (real.east) -- (sb.west);
\draw[arr] (real.east) -- (bo.west);
\draw[arr] (real.east) -- (mm.west);
\draw[arr] (real.north) |- (prvra.west);

\draw[arr] (cs.east) -- (prvfa.west);
\draw[arr] (sb.east) -- (pfvra.west);
\draw[arr] (bo.east) -- (pfvfa.west);
\draw[arr] (mm.east) -- (pmm.west);

\draw[arr] (prvra.east) -- (net.west |- prvra) -- (net.160);
\draw[arr] (prvfa.east) -- (net.west |- prvfa);
\draw[arr] (pfvra.east) -- (net.west |- pfvra);
\draw[arr] (pfvfa.east) -- (net.west |- pfvfa);
\draw[arr] (pmm.east)  -- (net.south west);

\draw[arr] (net.east) -- (lv.west);
\draw[arr] (net.east) -- (la.west);
\draw[arr] (net.east) -- (lc.west);

\draw[arr] (lv.east) -- (s0.west);
\draw[arr] (la.east) -- (s0.west);
\draw[arr] (lc.east) -- (s0.west);

\draw[arr, dashed] (s0.south) -- node[right]{warm start} (s1.north);

\end{tikzpicture}}%

  \caption{Quadrant-Aware Contrastive Pretraining (QACP). All four
  authenticity quadrants plus the real-but-mismatched class are constructed
  from pristine clips alone; three supervised-contrastive losses with
  different label partitions structure the three embedding spaces (Stage~0),
  after which the network warm-starts supervised multi-task fine-tuning
  (Stage~1).}
  \label{fig:qacp}
\end{figure}

\begin{table}[h]
\centering
\caption{QACP pseudo-classes constructed from pristine clips.}
\label{tab:qacp-classes}
\small
\begin{tabular}{lll}
\toprule
\textbf{Pseudo-class} & \textbf{Construction} & \textbf{Teaches} \\
\midrule
pseudo-RVRA & untouched real clip & anchor \\
pseudo-RVFA & vocoder copy-synthesis of the audio & audio authenticity cues \\
pseudo-FVRA & self-blended face video~\cite{shiohara2022sbi} & video authenticity cues \\
pseudo-FVFA & both transforms & consistent double-fake \\
pseudo-MISMATCH & real audio swapped from another clip & mismatch $\neq$ fake \\
\bottomrule
\end{tabular}
\end{table}

The MISMATCH class is the disentanglement forcing function: both streams remain
genuinely real while synchronization is broken, so the authenticity heads must
answer ``real'' while the consistency embedding flags disagreement. Three
supervised-contrastive (SupCon~\cite{khosla2020supcon}) losses are applied to
the three embedding spaces, each with a \emph{different} label partition of the
same batch:
\begin{equation}
\mathcal{L}_{\text{QACP}} =
  \mathcal{L}_{\text{SupCon}}(z_v;\, \text{video-auth}) +
  \mathcal{L}_{\text{SupCon}}(z_a;\, \text{audio-auth}) +
  \mathcal{L}_{\text{SupCon}}(z_c;\, \text{sync-state}).
\end{equation}
Because each space is supervised to be invariant to the other factors, the
factorization itself performs disentanglement.

\subsection{Deployment Architecture}
The complete framework is deployed as a two-tier system: a FastAPI inference
service hosting a lightweight DAVID-Net variant on Hugging Face Spaces, and a
Next.js web interface that presents per-modality verdicts, quadrant
attribution, temporal manipulation timelines, and the synchronization curve to
end users.

\section{Data Collection}
\label{sec:data}

\subsection{Selection of the Data Collection Method}
This research uses \emph{established benchmark datasets} rather than an
interview, survey, or systematic literature review as its primary data source.
This choice follows directly from the research questions: RQ1--RQ6 concern the
measurable behavior of a computational artifact on labeled audio-visual media,
which only curated, ground-truth-annotated corpora can answer. Interviews or
surveys measure human opinion and cannot supply the frame-level manipulation
labels required for training and for the localization task; benchmark corpora
additionally allow direct, reproducible comparison against prior work under
identical splits, which is the accepted evaluation standard in this field.

\subsection{Description of the Datasets}
Table~\ref{tab:datasets} lists the datasets used.
FakeAVCeleb~\cite{khalid2021fakeavceleb} provides native four-quadrant labels
and is the primary training set. AV-Deepfake1M~\cite{cai2023avdeepfake1m} and
LAV-DF~\cite{cai2022lavdf} drive the temporal-localization task.
DFDC~\cite{dolhansky2020dfdc} and KoDF are reserved exclusively for
cross-dataset evaluation. ASVspoof~\cite{wang2020asvspoof} strengthens the
audio branch. All datasets are obtained under their respective end-user license
agreements.

\begin{table}[h]
\centering
\caption{Datasets used in this study.}
\label{tab:datasets}
\small
\begin{tabular}{llll}
\toprule
\textbf{Dataset} & \textbf{Modality} & \textbf{Labels} & \textbf{Role} \\
\midrule
FakeAVCeleb   & A+V & 4-quadrant           & train + in-domain test \\
AV-Deepfake1M & A+V & clip + localization   & train + localization \\
LAV-DF        & A+V & localized forgeries   & localization \\
DFDC          & A+V & binary                & cross-dataset test only \\
KoDF          & A+V & binary                & cross-dataset test only \\
ASVspoof 2019/21 & A & bona fide / spoof    & audio-branch auxiliary \\
\bottomrule
\end{tabular}
\end{table}

\subsection{Preprocessing and Analysis Pipeline}
Video is decoded at 25\,fps; faces are detected, tracked, and aligned,
producing a $224\times224$ face crop and a $96\times96$ mouth crop. Audio is
resampled to 16\,kHz mono and loudness-normalized. All datasets are converted
to a unified manifest schema with per-modality labels, quadrant, manipulated
intervals, generator identity, and demographic metadata, making the analysis
reproducible from committed configuration files. Splits are subject-disjoint;
leave-one-generator-out (LOGO) splits hold out one manipulation method
entirely. The end-to-end data pipeline, from raw corpora to the deployed
service, is shown in Figure~\ref{fig:pipeline}.

\begin{figure}[t]
  \centering
  % ------------------------------------------------------------------
% Data collection / processing pipeline diagram (TikZ):
% raw corpora -> unified manifest -> preprocessing -> feature cache ->
% splits -> training / evaluation -> deployment.
% ------------------------------------------------------------------
\resizebox{\textwidth}{!}{%
\begin{tikzpicture}[
  font=\scriptsize,
  box/.style={draw, rounded corners=2pt, align=center, inner sep=5pt,
              minimum height=1.0cm, minimum width=2.3cm},
  data/.style={box, fill=gray!12},
  proc/.style={box, fill=blue!10},
  store/.style={box, fill=yellow!22},
  eval/.style={box, fill=green!14},
  dep/.style={box, fill=violet!12},
  arr/.style={-{Stealth[length=2.2mm]}, semithick}
]

% raw datasets
\node[data, minimum height=2.6cm, minimum width=2.6cm] (raw) at (0, 0)
  {Benchmark corpora\\[2pt]FakeAVCeleb\\AV-Deepfake1M\\LAV-DF, DFDC,\\KoDF, ASVspoof};

% manifest
\node[proc] (man) at (3.6, 0)
  {Unified manifest\\(per-modality labels,\\quadrant, segments,\\generator, meta)};

% preprocessing split into two
\node[proc, fill=blue!14] (vprep) at (7.2, 1.4)
  {Video: face detect,\\track, align;\\face + mouth crops};
\node[proc, fill=orange!18] (aprep) at (7.2, -1.4)
  {Audio: 16 kHz mono,\\loudness normalize};

% feature cache
\node[store] (cache) at (10.8, 0)
  {Feature cache\\(frozen VideoMAE /\\WavLM, one pass,\\local NVMe)};

% splits
\node[proc] (splits) at (14.2, 0)
  {Committed splits\\subject-disjoint /\\LOGO /\\cross-dataset};

% training & evaluation
\node[eval] (train) at (17.6, 1.4)
  {Stage 0 QACP +\\Stage 1 training\\(DGX Spark, BF16)};
\node[eval] (evaln) at (17.6, -1.4)
  {Evaluation: AUC, EER,\\quadrant, AP@IoU,\\ECE, robustness, fairness};

% deployment
\node[dep, minimum height=1.2cm] (deploy) at (21.2, 0)
  {Deployment:\\FastAPI (HF Space)\\+ Next.js UI};

% arrows
\draw[arr] (raw) -- (man);
\draw[arr] (man.east) -- (vprep.west);
\draw[arr] (man.east) -- (aprep.west);
\draw[arr] (vprep.east) -- (cache.west);
\draw[arr] (aprep.east) -- (cache.west);
\draw[arr] (cache) -- (splits);
\draw[arr] (splits.east) -- (train.west);
\draw[arr] (splits.east) -- (evaln.west);
\draw[arr] (train.south) -- (evaln.north);
\draw[arr] (train.east) -- (deploy.west);
\draw[arr] (evaln.east) -- (deploy.west);

\end{tikzpicture}}%

  \caption{The end-to-end data collection and processing pipeline. Encoder
  features are cached once (frozen backbones), making every subsequent
  training run inexpensive and every reported number reproducible from
  committed manifests and split files.}
  \label{fig:pipeline}
\end{figure}

Experiments are implemented in PyTorch and executed on an NVIDIA DGX Spark
system (GB10 Grace--Blackwell, 128\,GB unified memory) using BF16 precision and
a cached-feature training regime. Evaluation uses per-modality AUC, EER,
macro-F1, four-class quadrant accuracy, localization AP@IoU, expected
calibration error, robustness curves under compression and noise, and fairness
gaps across demographic subgroups; statistical rigor is ensured through
three-seed runs, bootstrap confidence intervals, and DeLong tests. Baselines
include Xception~\cite{rossler2019faceforensics},
LipForensics~\cite{haliassos2021lipforensics}, RawNet2~\cite{tak2021rawnet2},
AASIST~\cite{jung2022aasist}, MDS~\cite{chugh2020mds}, and
AVoiD-DF~\cite{yang2023avoiddf}, evaluated under identical splits.

\section{Professional Code of Ethical Issues}
\label{sec:ethics}

The ethical dimension of this research parses into four sub-problems, each with
a concrete mitigation.

\textbf{Consent and licensing of biometric data.} The datasets contain faces
and voices of real individuals. All corpora are used strictly under their
end-user license agreements, no raw media is redistributed, and only manifests,
split definitions, and model weights are released.

\textbf{Sampling bias and fairness.} Benchmark corpora over-represent certain
demographics and languages, which can bias detection accuracy. The study
audits performance gaps across gender, apparent skin tone, and language
subgroups, applies balanced sampling during training, and reports the residual
gaps transparently rather than concealing them.

\textbf{Transparency and research integrity.} All reported numbers correspond
to committed configuration files with fixed random seeds; negative and partial
results of ablations are reported alongside positive ones, guarding against
falsification, fabrication, and selective reporting. There is no conflict of
interest and no external funding influence on the findings.

\textbf{Dual-use and responsible deployment.} The system is strictly
defensive: no generation capability is included or released. The public
interface presents verdicts as calibrated probabilities with an explicit
``probabilistic, not proof'' disclaimer, user uploads are not retained beyond
the request, and the model card documents intended use and limitations.

\section{Economic Decision}
\label{sec:economic}

Two alternative strategies were compared for the experimental infrastructure.
\textbf{Alternative A (cloud GPU rental)}: training on rented A100-class
instances at current market rates, with an estimated total of
\todo{X}~GPU-hours across all runs, ablations, and seeds. \textbf{Alternative B
(owned DGX Spark workstation)}: a one-time hardware cost with negligible
marginal cost per experiment. Because the experimental plan requires 15--20
full training runs plus continuous development iteration, the owned-hardware
alternative amortizes favorably within a single semester and removes the risk
of budget-driven experiment truncation; it was therefore selected. The
cached-feature training regime further reduces compute cost by amortizing the
expensive encoder forward passes into a one-time preprocessing step.

For real-world deployment, the cost breakdown of the proposed solution is
deliberately minimal: the demonstration model is served from a free-tier
Hugging Face Space (CPU) with an optional GPU upgrade
(\todo{quote current rate}); the web interface is hosted on Vercel's free
tier; and the heavyweight model can be self-hosted on the same DGX Spark
workstation at an estimated \todo{X}~W average power draw. A platform-scale
adopter would instead face inference costs of approximately
\todo{estimate}~per thousand clips on commodity cloud GPUs, which the
lightweight distilled variant is specifically designed to reduce.
